
\FloatBarrier

\begin{center}
    {\Large \textbf{Supplementary Material for Déjà View}}
\end{center}



\section{Training Datasets}
\label{sec:training_datasets}
We train on a mixture of 29 publicly available datasets that span synthetic
renderings, indoor and outdoor real captures, multi-view object scans, and
driving footage. The corpus is highly imbalanced: per-dataset image counts
$N_i$ span more than three orders of magnitude (from $\sim$10\,k for
Spring to $\sim$11.8\,M for Aria Synthetic Environments), so naive
proportional sampling would let a handful of large datasets dominate
training, while uniform sampling would massively oversample the smallest
ones. In LLM literature, multilingual training faces the same imbalance across high- and
low-resource languages, which is addressed with temperature sampling
\citep{arivazhagan2019massively}. We adopt
the same recipe, treating each dataset as a ``language'' with token
budget $N_i$: the probability of drawing a training
example from dataset $i$ is set to
\[
    p_i \;=\; \frac{N_i^{\alpha}}{\sum_j N_j^{\alpha}},
    \qquad \alpha = 0.5,
\]
i.e.\ proportional to $\sqrt{N_i}$, which flattens the head of the
distribution while still favouring the larger, more diverse corpora.
The realised epoch share for each dataset is reported in
Table~\ref{tab:training_datasets}.

\FloatBarrier
\begin{table}[t]
    \centering
    \footnotesize
    \caption{Per-epoch training-mixture share for the 29 datasets, sorted
    by share. By construction, mix\,\% $= p_i \propto \sqrt{N_i}$ where
    $N_i$ is the total number of training images in dataset $i$.}
    \label{tab:training_datasets}
    \begingroup
\setlength{\tabcolsep}{3pt}
\begin{tabular}{@{}lr@{\hspace{0.8em}}lr@{}}
    \toprule
    \textbf{Dataset} & \textbf{mix \%} & \textbf{Dataset} & \textbf{mix \%} \\
    \midrule
    Aria Synth.\ Env.~\citep{avetisyan2024scenescript}                  & 13.77 & BEDLAM~\citep{black2023bedlam}                       & 1.99 \\
    WildRGB-D~\citep{xia2024wildrgbd}                                   &  9.32 & BlendedMVS~\citep{yao2020blendedmvs}                 & 1.35 \\
    TRELLIS~\citep{xiang2025trellis,deitke2023objaverse}                &  9.04 & MegaDepth~\citep{li2018megadepth}                    & 1.16 \\
    CO3D~\citep{reizenstein2021co3d}                                    &  8.84 & Replica~\citep{straub2019replica}                    & 0.95 \\
    Taskonomy~\citep{zamir2018taskonomy}                                &  7.26 & Virtual KITTI~2~\citep{cabon2020vkitti2}             & 0.83 \\
    ScanNet~\citep{dai2017scannet}                                      &  6.23 & Hypersim~\citep{roberts2021hypersim}                 & 0.69 \\
    DL3DV~\citep{ling2024dl3dv}                                         &  5.84 & Kubric~\citep{greff2022kubric}                       & 0.60 \\
    Cubify Any.~\citep{lazarow2025cubifyanything,baruch2021arkitscenes} &  5.73 & GTA-SfM~\citep{wang2020gtasfm}                       & 0.52 \\
    TartanAir~V2~\citep{wang2020tartanair}                              &  4.78 & MatrixCity~\citep{li2023matrixcity}                  & 0.49 \\
    Parallel Dom.\ 4D~\citep{vanhoorick2024gcd}                         &  4.75 & MPSD~\citep{lopez2020mpsd}                           & 0.49 \\
    ScanNet++~\citep{yeshwanth2023scannetpp}                            &  3.96 & Mapillary Metr.~\citep{mapillary2021metropolis}      & 0.47 \\
    Waymo~\citep{sun2020waymo}                                          &  2.76 & UnrealStereo4K~\citep{tosi2021smdnets}               & 0.44 \\
    Mid-Air~\citep{fonder2019midair}                                    &  2.61 & MVS-Synth~\citep{huang2018deepmvs}                   & 0.44 \\
    Dynamic Replica~\citep{karaev2023dynamicstereo}                     &  2.16 & Spring~\citep{mehl2023spring}                        & 0.40 \\
    TartanAir~\citep{wang2020tartanair}                                 &  2.13 &                                                      &      \\
    \bottomrule
\end{tabular}
\endgroup

\end{table}


\section{Two-Stage Depth Training}
\label{sec:suppl_two_stage}


The linear pixel-shuffle head maps patch tokens to per-pixel depth via pixel shuffling. For ray directions, this is acceptable as patch-border gradients are within $\sim$10\% of intra-patch values. For depth, the same gradients are roughly an order of magnitude larger, producing visible block artifacts at patch boundaries (\cref{fig:suppl_grid_artifacts}).

We address this with two-stage training (\cref{sec:method_training}). The first stage trains the model end-to-end with a linear pixel-shuffle depth head and plain $\ell_2$ depth loss. Training the final pipeline configuration (convolutional head with confidence loss) end-to-end instead yields worse metrics across our benchmarks (\cref{tab:suppl_two_stage_ablation}). The second stage swaps it in, freezes the rest of the network, and finetunes the depth decoder with a confidence-weighted loss. Its convolutions smooth across patch boundaries, eliminating the block pattern (\cref{fig:suppl_grid_artifacts}) and yielding a slight improvement in Inlier Ratio (\cref{tab:suppl_full_recipe_two_stage}). As a result, we also obtain a depth confidence channel that can be used downstream to filter regions with uncertain reconstruction.

\begin{table}[t]
    \caption{%
        \textbf{End-to-end vs.\ two-stage training.}
        Training the final pipeline configuration end-to-end (top, convolutional head with confidence loss) consistently underperforms the first stage of our recipe (bottom, linear head with $\ell_2$ depth loss) on every metric.
    }
    \label{tab:suppl_two_stage_ablation}
    \footnotesize   
    \setlength{\tabcolsep}{6pt}

    \noindent\makebox[\linewidth][c]{%
    \begin{tabular}{l rrrr}
        \toprule
        \textbf{Variant}
            & \textbf{AUC@$3^{\circ}$}~$\uparrow$
            & \textbf{AUC@$30^{\circ}$}~$\uparrow$
            & \textbf{IR}~$\uparrow$
            & \textbf{AbsRel}~$\downarrow$ \\
        \midrule
        Convolutional head + confidence loss (end-to-end) & 23.0 & 77.1 & 56.5 & 0.152 \\
        \midrule
        \textbf{Linear head + $\ell_2$ (our stage 1)}     & \textbf{31.0} & \textbf{80.6} & \textbf{59.2} & \textbf{0.125} \\
        \bottomrule
    \end{tabular}}
\end{table}

\begin{table}[t]
    \caption{%
        \textbf{Stage 2 finetuning.} Finetuning our stage 1 model (linear head + $\ell_2$ depth loss) with the convolutional head and confidence-weighted loss (stage 2) yields a small, consistent gain in Inlier Ratio while leaving the pose metrics unchanged.
    }
    \label{tab:suppl_full_recipe_two_stage}
    \footnotesize
    \setlength{\tabcolsep}{6pt}

    \noindent\makebox[\linewidth][c]{%
    \begin{tabular}{l rrrr}
        \toprule
        \textbf{Variant}
            & \textbf{AUC@$3^{\circ}$}~$\uparrow$
            & \textbf{AUC@$30^{\circ}$}~$\uparrow$
            & \textbf{IR}~$\uparrow$
            & \textbf{AbsRel}~$\downarrow$ \\
        \midrule
        Linear head + $\ell_2$ (stage 1, full recipe) & \textbf{56.8} & \textbf{91.8} & 79.8 & \textbf{0.031} \\
        \midrule
        \textbf{Conv. head + conf. loss (stage 2 finetune)} & \textbf{56.8} & \textbf{91.8} & \textbf{80.3} & \textbf{0.031} \\
        \bottomrule
    \end{tabular}}
\end{table}



\begin{figure}[t]
    \centering
    \begin{tabular}{@{}c@{\hspace{0.02\linewidth}}c@{\hspace{0.02\linewidth}}c@{}}
        \textbf{Input} & \textbf{Linear Head} & \textbf{Convolutional Head} \\[2pt]
        \includegraphics[width=0.32\linewidth]{figures/block_artifacts_rgb.png} &
        \includegraphics[width=0.32\linewidth]{figures/block_artifacts_stage1.png} &
        \includegraphics[width=0.32\linewidth]{figures/block_artifacts_stage2.png}
    \end{tabular}
    \caption{%
        \textbf{Block artifacts.}
        Depth predicted by the first-stage linear head \textbf{(center)} shows visible block artifacts aligned with the DINOv2 patch grid. The second-stage finetune with the convolutional head \textbf{(right)} eliminates them and improves Inlier Ratio.
    }
    \label{fig:suppl_grid_artifacts}
\end{figure}

\section{Emergent Iterative Correspondence Search}
\label{sec:emergent_correspondence}

\begin{figure*}[t]
    \centering
    \includegraphics[width=\linewidth]{figures/feature_match.pdf}
    \caption{
        \textbf{Emergent iterative correspondence search.}
        For two example query patches (\textcolor{green}{green square}) 
        we visualize the head-averaged global self-attention
        sub-block weights induced by that query at each iteration of
        the loop, overlaid on the corresponding target view.
        Iterations advance left-to-right, the attention starts diffuse
        and progressively concentrates on the corresponding counterpart of
        the queried patch, despite the absence of any explicit feature matching
        supervision during training.
    }
    \label{fig:feature_correspondence}
\end{figure*}


We probe the global attention sub-block of our recurrent layer to investigate how its attention pattern evolves across iterations. For a query patch in view~0, at each iteration $t$ we read the per-head queries and keys $Q^{(t)}_h, K^{(t)}_h$ post $q$ -/ $k$ -LayerNorm and visualize the head-averaged scaled-dot-product attention weights
\[
    \bar a^{(t)} \;=\;
    \tfrac{1}{H} \sum_{h=1}^{H} \operatorname{softmax}\!\left(
        q^{(t)}_h\, K^{(t)\top}_h \big/ \sqrt{d_h}
    \right),
\]
sliced to the patch tokens of each target view and shown as a $H_p \times W_p$ heatmap. 
\cref{fig:feature_correspondence} illustrates that the attention starts diffuse and progressively concentrates on the corresponding counterpart of the queried patch, correctly resolving symmetries and attending to all geometrically equivalent patches. This suggests that the recurrent loop implements an emergent iterative correspondence search, despite the model being supervised only on 3D reconstruction and pose estimation losses.


\section{Scaling Beyond $K_{\max}$}
\label{sec:beyond_kmax}

\begin{figure*}[t]
    \centering
    \includegraphics[width=\linewidth]{figures/beyond_kmax_figure.pdf}
    \caption{
        \textbf{Step-count extrapolation.} Test-time metrics as a function
        of the inference step count $K_{\inf}$. Performance peaks
        at the maximum trained budget, then degrades when moving far outside the trained range.
    }
    \label{fig:beyond_kmax}
\end{figure*}


DéjàView is trained with an iteration count sampled per-batch as $K \sim \mathrm{Beta}(2,1)$
scaled to $[K_{\min}, K_{\max}]$ with $K_{\max}=16$. We sweep $K_{\inf}$ at test time beyond $K_{\max}$ and observe that Pose~AUC@$3^\circ$ peaks
near the trained budget, then starts to degrade after. Pose~AUC@$30^\circ$ and Pointmap Rel.~L2
remain stable up to approximately $K_{\inf}=30$ but eventually collapse. \cref{fig:beyond_kmax} shows that model iterations cannot be pushed arbitrarily far. 

Our analysis shows that this occurs because some feature channels grow unbounded as we scale beyond $K_{\max}$. The mechanism is already visible inside the trained range (\cref{fig:refinement}b): while cosine similarity to
$\mathbf{z}_{K_{\max}}$ saturates near $1$ and the relative feature update decays, the state norm $\|\mathbf{z}_k\|_2$ keeps growing monotonically through $K_{\max}$ after a short initial contraction. Iterating beyond $K_{\max}$ simply extrapolates this persistent drift: a handful of channels grow without bounds, producing the observed collapse.

\section{Scaling Below $K_{\max}$}
A sub-$K_{\max}$ compute allows two strategies: a full $K_{\inf}$-step forward
with uniform time interval conditioning calibrated for that budget, or early-stopping a
$K_{\max}$-step rollout by reading $\mathbf{z}_k$ at $k=K_{\inf}$.
\cref{fig:kinf_intermediate} shows the full pass strictly dominates
early-stopping for every $k<K_{\max}$, with the largest gap at the smallest
budget. At $K_{\inf}=8$, Pose~AUC@$3^\circ$ rises from $0.31$ to $0.44$
($\sim\!43\%$ relative improvement).


\begin{figure*}[t]
    \centering
    \includegraphics[width=\linewidth]{figures/overall_metrics_comparison.pdf}
    \caption{
        \textbf{Decoding intermediate results vs. using lower $K_{\inf}$.}
        We compare using $K_{\inf} < K_{\max}$ steps explicitly to decoding $\mathbf{z}_k$ at intermediate iterations $k$ when $K_{\inf} = K_{\max}$, and show that explicitly conditioning the model on fewer iterations degrades performance less than early-stopping. 
    }
    \label{fig:kinf_intermediate}
\end{figure*}



\section{Baseline Configurations}
\label{sec:suppl_baselines}

We run all baselines through our evaluation framework using their official code releases and checkpoints, and apply the same Sim(3) alignment to predicted pointmaps before computing metrics (\cref{sec:sota}). For Rel.~L2 and IR we use each method's primary 3D output: depth-unprojected pointmaps for VGGT, DA3, and our model, the direct point-head output for Pi3 and MapAnything, and the SGA-optimized dense pointmap for MASt3R and MASt3R-SfM.

\paragraph{VGGT~\citep{vggt}.}
The official \texttt{facebook/VGGT-1B} checkpoint at $518$-pixel longest edge with patch size $14$.

\paragraph{VGGT-$\Omega$~\citep{vggt_omega}.}
The official \texttt{facebook/VGGT-Omega} 1B checkpoint (\texttt{vggt\_omega\_1b\_512.pt}, without text alignment) at $512$-pixel longest edge with patch size $16$. We decode the released 9D camera encoding (translation, quaternion, FoV$_h$, FoV$_w$) via the official \texttt{encoding\_to\_camera} utility and use depth-unprojected pointmaps as the primary 3D output, matching the convention used for VGGT and DA3.

\paragraph{Pi3~\citep{wang2025pi}.}
The official \texttt{yyfz233/Pi3} checkpoint at $518$-pixel longest edge with patch size $14$. %

\paragraph{MapAnything~\citep{keetha2026mapanything}.}
The official \texttt{facebook/map-anything} v1.1 checkpoint at $518$-pixel longest edge with patch size $14$.

\paragraph{Depth Anything 3 (DA3)~\citep{lin2025depth}.}
The official v1.1 checkpoints at two backbone scales (DA3-L: \texttt{depth-anything/DA3-LARGE}, ViT-L, $356$M params; DA3-G: \texttt{depth-anything/DA3-GIANT}, ViT-G, $1.2$B params), both at $504$-pixel longest edge with patch size $14$. We decode camera pose from predicted rays.

\paragraph{MASt3R~\citep{leroy2025grounding}.}
The official \texttt{MASt3R\_ViTLarge\_BaseDecoder\allowbreak\_512\_catmlpdpt\_metric} (metric-scale) checkpoint at $512$-pixel longest edge with patch size $16$, using DUSt3R-style image normalization (mean and std $0.5$). Pair selection is adaptive on scene size: \texttt{complete} (all pairs) for scenes with $N \le 8$ views, and \texttt{swin-5} (sliding window of $5$) otherwise. The efficiency measurement at $N{=}24$ uses the swin-5 branch ($120$ pairs). The sparse global alignment uses the published defaults: $300$ iterations of coarse alignment at learning rate $0.07$ followed by $300$ iterations of refinement at $0.01$, with per-pixel depth optimization enabled (\texttt{optim\_level=refine+depth}) and matching-confidence threshold $5.0$. Camera intrinsics are shared across views for the single-camera benchmarks (7-Scenes, ScanNet++, nuScenes, DTU) and estimated per-view otherwise (ETH3D).

\paragraph{MASt3R-SfM~\citep{duisterhof2024mast3rsfm}.}
The same MASt3R checkpoint paired with the official training-free retrieval model (\texttt{MASt3R\_ViTLarge\_BaseDecoder\_512\_catmlpdpt\_metric\_retrieval\_trainingfree}). Scene graphs are built via top-$20$ retrieval anchors with top-$10$ retrieved neighbors per anchor (\texttt{retrieval-20-10}), yielding $273$ pairs at $N{=}24$. The sparse global alignment uses the same hyperparameters as MASt3R, including the same per-dataset shared-intrinsics policy.


\section{Societal Impact}
\label{sec:suppl_societal_impact}
\methodname{} reconstructs 3D geometry and camera poses from images. While this capability is not new, our work brings strong reconstruction quality within reach at a smaller scale than recent feed-forward baselines, lowering the overall cost of deployment. Because the model outputs geometry rather than photorealistic imagery, the direct risk of deceptive media generation is lower than for image or video synthesis models, though deployments that combine reconstruction with generative rendering should still consider provenance signals such as watermarking~\citep{wen2023treering}. From an environmental perspective, our ViT-B model is trained on 128 H100 GPUs, comparable to other recent feed-forward reconstruction methods. At inference, however, it operates with roughly an order of magnitude fewer parameters and a memory footprint of under 5 GiB at $24$ input views, reducing the per-query resource cost of deployment relative to the larger baselines we evaluate.

\begin{figure*}[t]
    \centering
    \includegraphics[width=\linewidth]{figures/qualitatives2.pdf}
    \caption{
        \textbf{Qualitative results.}
        Predicted point clouds and cameras for in-the-wild captures.
    }
    \label{fig:qualitative2}
    \vspace{-4mm}
\end{figure*}







